%------------------------------------------------------------------------------
% \PART 2
%------------------------------------------------------------------------------
\chapter[Анализ технического задания]{АНАЛИЗ ТЕХНИЧЕСКОГО ЗАДАНИЯ}

%------------------------------------------------------------------------------
% \PART 2.1 Text
%------------------------------------------------------------------------------
\section{Постановка задачи}

\hspace*{12.5 mm}В данной работе основное внимание уделяется задаче 
распознавания изображений рукописных цифр с использованием нейронных сетей, 
оптимизированных для аппаратной реализации на ПЛИС (FPGA).

Целью дипломной работы является проектирование и реализация компактной 
нейронной сети прямого распространения, способной обрабатывать изображения 
размером 28$\times$28 пикселей, представляющих рукописные цифры из 
стандартного набора данных MNIST. Основной задачей является не только 
достижение высокой точности классификации, но и обеспечение совместимости с 
аппаратной реализацией, особенно в условиях ограниченных вычислительных и 
энергетических ресурсов.

В качестве основы для нейросетевой архитектуры используется метод Learned 2D 
Separable Transform, предложенный в работе~\cite{LST2D}. Преимущество данного 
подхода заключается в использовании разделяемых по строкам и столбцам операций, 
что позволяет существенно снизить количество параметров модели по сравнению с 
традиционными полносвязными слоями. LST-архитектура даёт возможность эффективно 
использовать параллелизм, характерный для FPGA, и обеспечивает высокую скорость 
обработки без значительных потерь в точности.

Для аппаратной реализации модели выбрана микросхема XC7Z010CLG400-2, входящая в 
состав популярной отладочной платы Zybo Z7. Это решение основано на архитектуре 
Xilinx Zynq-7000, сочетающей возможности ПЛИС и встроенного процессора ARM, что 
делает его удобным для интеграции программно-аппаратных систем.

Вид сверху отладочной плата Zybo Z7 на базе Xilinx Zynq-7000, с указанием 
основных интерфейсов и компонентов: USB, HDMI, DDR, питание, PMOD-разъемы и т.д. 
представлен на рисунке~\ref{fig:z7}~\cite{ZYBO}.

\insertfigure{z7}{pic/zybo-z7.png}{ПЛИС Zybo Z7}{12cm}

Внутренняя структура Zybo Z7 представлена на рисунке~\ref{fig:z7-in}.

\insertfigure{z7-in}{pic/zybo-in.jpg}{Структура ПЛИС Zybo Z7}{16.0cm}

На схеме продемонстрированы основные блоки SoC Zynq-7000, включая: процессорную 
систему (PS), программируемую логику (PL), AXI-интерфейсы, контроллеры памяти и 
периферийные модули. В рамках настоящего проекта будет использоваться 
взаимодействие между PL и PS через AXI-шину

Для взаимодействия с моделью и удобства проведения тестирования будет 
использоваться программная платформа PYNQ, позволяющая запускать и отлаживать 
алгоритмы с помощью Jupyter Notebook и Python API. Таким образом, будет 
реализован полный цикл: от проектирования и обучения нейросети – до её 
верификации и аппаратной имплементации.

Постановка задачи и предъявленные требования представлены в графическом 
материале ГУИР 431289.001 ПЛ.

%------------------------------------------------------------------------------
% \PART 2.2 Text
%------------------------------------------------------------------------------
\section{Анализ требований к нейронной сети}
\hspace*{12.5 mm}Выбор архитектуры нейронной сети требует учета как точности 
модели, так и её аппаратной эффективности. В условиях ограниченных ресурсов 
FPGA особое внимание уделяется упрощению структуры и снижению объема 
вычислений. Основные требования к нейросети включают:

  1 Модель должна обеспечивать точность не ниже 97\% на тестовой выборке MNIST, 
что соответствует уровню, достигаемому современными методами машинного 
обучения. Для достижения данной цели необходимо правильно подобрать 
архитектуру, функцию потерь, функцию активации и метод оптимизации.

  2 Для реализации на аппаратном уровне необходимо минимизировать количество 
параметров модели, сохраняя при этом высокую точность. Это особенно важно для 
устройств с ограниченными ресурсами, таких как встраиваемые системы.

  3 Память в ПЛИС системах имеет жёсткие ограничения. Модель должна быть 
оптимизирована таким образом, чтобы не только сама структура занимала как можно 
меньше памяти, но и чтобы объем оперативных данных, хранимых во время выполнения 
активации, был минимален.

  4 Выходной слой сети должен использовать softmax для получения вероятностного 
распределения по классам. В скрытых слоях могут использоваться функции 
активации, оптимальные для аппаратной реализации, например tanh. 

%------------------------------------------------------------------------------
% \PART 2.3 Text
%------------------------------------------------------------------------------
\section{Анализ требований к аппаратной реализации}
\hspace*{12.5 mm}Эффективность работы нейронной сети на аппаратном уровне 
зависит от грамотного распределения вычислительных ресурсов, схемотехнической 
оптимизации и способности к параллелизму. 

Основные аппаратные требования и ограничения при реализации модели на FPGA 
включают:

  1 В архитектуре FPGA присутствуют специализированные блоки: логические ячейки 
(LUT), регистры, блоки умножения (DSP), встроенная память (BRAM). Для успешной 
реализации модели необходимо уйти от операций с плавающей запятой, заменяя их 
на фиксированную арифметику, а также эффективно задействовать ресурсы DSP и 
BRAM. В FPGA важно использовать внутренний параллелизм.

  2 Так как встраиваемые системы часто работают от автономных источников 
питания, архитектура должна потреблять как можно меньше энергии. Этого можно 
достичь путём сокращения количества операций, понижения разрядности чисел, а 
также выбора энергоэффективных функций активации и структур слоёв.

  3 Архитектура должна быть масштабируемой. 
Кроме того, должна быть возможность модификации модели для адаптации под новые
задачи.

  4 Реализация должна быть совместима с использованием Python-интерфейсов и 
поддерживать управление через PYNQ API. Это необходимо для быстрой верификации 
модели, автоматизации тестирования и сбора метрик производительности.

Основной задачей является выбор оптимальной архитектуры сети и методов 
вычислений, которые обеспечат баланс между производительностью и потреблением 
ресурсов.

%------------------------------------------------------------------------------
% \PART 2.4 Text
%------------------------------------------------------------------------------
\section{Выбор и обоснование метода решения задачи}
\hspace*{12.5 mm}Выбор архитектуры нейросети — один из ключевых аспектов в 
реализации высокоэффективного классификатора. Для решения задачи распознавания 
рукописных цифр рассматриваются различные подходы, каждый из которых имеет свои преимущества и 
ограничения.

Полносвязные нейронные сети  – обладают простой структурой и легко реализуются, 
однако требуют большого количества параметров. Это делает их менее подходящими 
для аппаратной реализации, так как объем необходимых вычислений и памяти растет
пропорционально числу нейронов в слоях.

Сверточные нейронные сети – обеспечивают высокую точность за счет эффективного 
извлечения признаков из изображений. Они используют свертки, позволяющие 
минимизировать количество параметров по сравнению с полносвязными сетями. 

Рекуррентные нейронные сети – могут использоваться для последовательной 
обработки данных, однако их применение к изображениям ограничено. Такие 
сети лучше подходят для обработки временных рядов и речи, но могут быть 
адаптированы к изображениям.

Трансформеры – показывают высокую производительность в задачах обработки 
естественного языка и компьютерного зрения, но требуют значительных 
вычислительных ресурсов. Несмотря на их преимущества, высокая сложность и 
потребление памяти делают их менее подходящими для работы на FPGA.\@

Гибридные архитектуры – могут комбинировать различные методы, такие как CNN и 
рекуррентные слои, что позволяет достичь оптимального баланса между точностью и
вычислительными затратами.

LST2D архитектуры — выступают в качестве адьтернативы архитектуре полносвязных 
слоев и позволяют использовать минимальные вычислительные ресурсы.

На основании требований к аппаратной реализации и соотношения сложности 
структуры доступных архитектур, предлагается использовать обучаемое двумерное 
разделимое преобразование, которое можно рассматривать как новый тип 
вычислительного слоя для построения различных архитектур нейронных сетей\cite{LST2D}. 

Данная структура обеспечивает высокую точность классификации изображений 
благодаря эффективному представлению пространственных зависимостей и 
позволяет сократить вычислительные затраты за счет оптимизированной структуры 
слоев.

Двумерный входной массив сначала преобразуется по строкам с помощью первого 
линейного слоя, а затем – по столбцам с использованием второго слоя с 
отдельными весами. Это позволяет:

  – сократить число параметров до $O(n^2)$ вместо $O(2n)$;

  – повысить скорость выполнения;

  – минимизировать потребление памяти.

Для реализации поставленных задач необходимо разработать структурную и 
функциональную схемы устройства. Описать алгоритм работы с учетом всех 
вышеперечисленных функций. По алгоритму описать модель на HDL языках.